% Fifth Philippic --- headnote
% philippic-5, 1 January 43 BC, Rome

\textit{The first speech of the consular year of Pansa and
Hirtius, delivered at the inaugural meeting of the Senate in
the temple of Concord on the Kalends of January, 43 BC. The
business of the day was the state of the war against Antony,
who since the autumn had been besieging Decimus Brutus at
Mutina; the question on the floor was a proposal --- carried,
it seems, by Quintus Fufius Calenus, the consular asked his
opinion first --- to send a legation to Antony rather than
proceeding at once to war. Cicero's speech is the great
rebuttal of that proposal and, with it, of every counsel of
delay. It is the longest speech in the cycle so far, and the
most legislative in shape: not the popular indignation of
Philippic 4 nor the formal motion of Philippic 3, but a
sustained debate-speech that begins with the case against the
legation, runs through Antony's whole record of public
violence, and closes with a series of drafted senatus consulta
honouring those who have already taken the field against him.}

\textit{The argument falls into two unequal halves. \textsection
1--34 are the case against the embassy and for war. \textsection
1--6 set the stake: a legation would soften both Roman opinion
and the resolve of the legions and municipalities that are
already in arms; it would also be incompatible with the honours
the Senate has just decreed (thirteen days before, in Philippic
3) to those who took up arms against Antony --- you cannot
both reward the Martian and Fourth Legions for abandoning him
and treat him as a citizen worth negotiating with. \textsection
7--16 turn to Antony's record as consul: the agrarian law
forced through with Lucius Antonius in defiance of Jove
thundering and the lex Caecilia et Didia; the embezzlement of
seven hundred million sesterces from the treasury of Ops; the
sale of decrees, kingdoms, citizenships, immunities; the
falsification of senatus consulta; and above all the third
panel of judges --- Greeks, gamblers, exiles, dancers, the
chorus of his revels --- which is held up at length for
ridicule. \textsection 17--22 cover the armed gangs around the
temple of Concord, the threat to demolish Cicero's house on
the Kalends of September, the speech of nineteenth September
in his absence, the centurions cut down before Antony and his
wife at Brundisium. \textsection 23--25 take the flight from
Rome to the siege of Mutina, and set the Hannibal comparison
(\textit{ergo Hannibal hostis, civis Antonius?}, \textsection
25): an enemy at the gates of a Roman colony, and to such a
man the proposal is to send legates. \textsection 26--34 are
the practical case against the embassy --- it will quench the
fire, it will breed doubt, doubt will paralyse the levy ---
and close with Cicero's own motion: \textit{tumultus} decreed,
\textit{iustitium} edicted, \textit{saga} taken up, a
universal levy held throughout Italy except Gaul, the
commonwealth committed to the consuls with the formula \textit{
ne quid res publica detrimenti accipiat}, and an amnesty for
any of Antony's soldiers who desert him before the Kalends
of February.}

\textit{\textsection 35--53 are the second half: the honours.
Cicero proposes formal senatus consulta for Decimus Brutus
(praise for holding Gaul against Antony, \textsection 36),
for Lepidus (a gilded equestrian statue on the Rostra in
return for restoring Sextus Pompey to the citizenship,
\textsection 41), for Gaius Caesar --- that is, Octavian, who
is throughout this speech \textit{C.~Caesar} --- propraetorian
\textit{imperium}, a seat in the Senate with consular rank,
and the right to stand for office as if he had held the
quaestorship the year before (\textsection 46); for Lucius
Egnatuleius, who brought the Fourth Legion across, the right
to seek office three years before the legal age (\textsection
52); and finally for the veterans of the Martian and Fourth
Legions a sweeping package of exemptions, land, and bonuses
(\textsection 53). The Caesar passage (\textsection 42--51) is
the longest single block: a defence of the unprecedented
honours against the foreseeable objection that the young man
is too young, with Alexander, Scipio Africanus, and Flamininus
brought in as precedents, and a personal guarantee of Caesar's
future conduct (\textit{promitto, recipio, spondeo}, \textsection
51) which the next eighteen months would expose. Almost all of
this carried, except the demand for \textit{tumultus} and \textit{
saga sumi}: the Senate compromised and did send the legation
(Sulpicius, Piso, Philippus), which is the business of
Philippic 6 and 7.}

\textit{The register is the high deliberative one. There is no
direct address to the people --- this is throughout \textit{
patres conscripti}, and the rhythms are those of the consilium
publicum rather than the contio. The polemic is sharper than
in Philippic 3 (Antony is now \textit{scelerati gladiatoris
amentia}, \textsection 32, and the great extended caricature
of the third panel of judges, with Cydas of Crete and Lysiades
the Athenian named, is among the most savagely funny set-
pieces in the cycle), but it is set within a thoroughly
legislative argument: the speech ends not in a peroration but
in the consul's formula \textit{dixi ad ea omnia, consules, de
quibus rettulistis}, the senatorial equivalent of resting one's
case.}
